% =============================================================================
% The GCHP CO2 4D-Var software: architecture, data flow and workflow
% Companion to:
%   monthly_4dvar_science.tex   - scientific formulation and results
%   monthly_sigma_gradient.tex  - derivation of the monthly adjoint gradient
% =============================================================================
\documentclass[11pt]{article}

% Latin Modern ships Type 1 outlines; plain Computer Modern needs MetaFont (mf),
% absent on the NAS hosts.  Do not add T1 fontenc or the bm package.
\usepackage{lmodern}
\usepackage{amsmath,amssymb}
\usepackage{geometry}
\usepackage{booktabs}
\usepackage{longtable}
\usepackage{hyperref}
\usepackage{xcolor}
\usepackage{enumitem}

\geometry{margin=1in}
\hypersetup{colorlinks=true, linkcolor=blue!50!black, urlcolor=blue!50!black}

% Compact monospace helper
\newcommand{\code}[1]{\texttt{#1}}

\title{\textbf{The GCHP CO$_2$ 4D-Var System:\\
Software Architecture and Workflow}}
\author{}
\date{\today}

\begin{document}
\maketitle

\begin{abstract}
This document describes the software that implements the monthly 4D-Var CO$_2$
flux assimilation: the directory layout, the programs and their interfaces, the
per-iteration workflow, the on-disk data formats, the GCHP configuration
(including the forward/adjoint asymmetry that makes the adjoint gradient correct),
the monthly-specific plumbing, the variable-window parameterisation, the
validation tooling, and the operational concerns (job chaining, walltime budgets,
checkpoint disk usage). It is a reference for running, debugging, and extending
the system. The scientific formulation is in \code{monthly\_4dvar\_science.tex}
and the gradient derivation in \code{monthly\_sigma\_gradient.tex}.
\end{abstract}

\tableofcontents

% =============================================================================
\section{Overview}
% =============================================================================

The system is a driver script that orchestrates repeated forward and adjoint
GCHP integrations and a quasi-Newton update, plus a set of Python programs that
handle the control variable, the observation forcing, the gradient, and the
diagnostics. One outer iteration of the assimilation consists of four phases:

\begin{enumerate}[nosep]
  \item \textbf{sigma + forward}: write the monthly scaling fields from the
        current state, run the forward GCHP integration, and move the adjoint
        checkpoints into place;
  \item \textbf{forcing}: compute the observation cost $J_{\mathrm{obs}}$ and the
        adjoint forcing from the simulated satellite tracks;
  \item \textbf{adjoint}: run the backward GCHP integration to produce the surface
        adjoint accumulator, then free the consumed checkpoints;
  \item \textbf{update}: form the gradient, take one L-BFGS-B step, snapshot the
        state, and generate diagnostic plots.
\end{enumerate}

Two variants coexist and are fully isolated from one another: a single-window
(``one-month'') optimiser and the monthly optimiser. They share the Python
components wherever possible; the monthly optimiser differs in the control
variable (twelve fields), the forward input plumbing (a time-varying monthly
field), and the gradient (boundary differencing).

% =============================================================================
\section{Directory layout}
% =============================================================================

Paths below are relative to the project root
\code{gchp\_14.5.3\_adjoint\_surfaceF/}. The scripts live in \code{gchp\_4Dvar/};
the run trees are siblings.

\begin{center}
\small
\begin{longtable}{@{}p{0.42\textwidth}p{0.52\textwidth}@{}}
\toprule
Path & Role \\
\midrule
\endhead
\code{gchp\_4Dvar/} & scripts, control-file templates, this documentation \\
\code{gchp\_4Dvar/control\_files/forward\_osse\_monthly/} & forward GCHP config (monthly) \\
\code{gchp\_4Dvar/control\_files/adjoint\_osse\_monthly/} & adjoint GCHP config (monthly) \\
\code{gchp\_4Dvar/control\_files/forward\_osse\_monthly\_test/} & forward config for the plumbing test \\
\code{4dvar\_design\_osse\_monthly/} & monthly run tree (work dir) \\
\quad \code{forward\_run/} & forward executable, restarts, \code{OutputDir} \\
\quad \code{adjoint/} & adjoint executable, restarts, \code{OutputDir}, checkpoints \\
\quad \code{sigma\_monthly/YYYY/MM.nc4} & monthly control fields \\
\quad \code{forcing\_files\_osse/} & adjoint forcing read by the adjoint run \\
\quad \code{4dvar\_output\_osse/} & per-iteration state snapshots and forcing archives \\
\quad \code{plots\_osse/<stamp>/} & diagnostic figures, one subdir per job \\
\code{4dvar\_design\_osse/} & the single-window run tree (one-month variant) \\
\bottomrule
\end{longtable}
\end{center}

% =============================================================================
\section{Programs and interfaces}
% =============================================================================

\subsection{Driver: \code{4dvar\_optimizer.osse.monthly.run}}

A PBS batch script that runs the outer loop. Responsibilities: source the GCHP
and Python environments; sync control-file templates into the run tree; patch the
run timing from the window; run the four phases per iteration; checkpoint progress
and chain jobs when the walltime budget is spent. It is fully parameterised
through environment variables (Section~\ref{sec:running}).

\subsection{Control variable: \code{write\_sigma\_monthly.py}}

Writes the monthly scaling fields from the current state vector. Each month is a
single-record NetCDF file \code{sigma\_monthly/YYYY/MM.nc4} holding the variable
\code{Sigma\_CO2} on the $72\times46$ lat--lon grid, stamped at the month start.
The twelve-single-file layout (rather than one twelve-record file) mirrors an
existing input dataset known to work with the GCHP input component under a
monthly refresh template. Key options:
\begin{itemize}[nosep]
  \item \code{--state-file}: read the control vector \code{x\_next}; absent, write
        the prior $\sigma\equiv1$;
  \item \code{--n-months} (default 12) and \code{--start-month} (default 1):
        number of monthly fields and the calendar month of the first, with year
        rollover---this is what makes a shorter window possible;
  \item \code{--const-per-month}: bypass the state file with explicit per-month
        constants (used only by the plumbing test).
\end{itemize}
The control vector is $(\code{n\_months}\times\code{nlat}\times\code{nlon})$
flattened in C order with month index $0$ = the start month.

\subsection{Observation cost and forcing: \code{co2\_adjoint\_forcing\_osse.py}}

Given the simulated satellite-track files and the observations, computes
$J_{\mathrm{obs}}$ (equation~\eqref{eq:jobs} of the science note) and writes the
per-timestep adjoint forcing that the backward run injects. It accepts \emph{all}
monthly satellite-track files for the window as separate arguments, processes them
one at a time (a year of tracks is never concatenated in memory), and sums $J$
across files. Before the heavy observation matching it performs cheap, loud
pre-checks on the file time coordinates:
\begin{itemize}[nosep]
  \item the number of files must equal the number of months spanned by the window;
  \item every window month must be covered by exactly one file (a month in two
        files signals a stale track from a previous run; a month in none signals
        silently dropped observations---both are hard errors);
  \item stale forcing outputs from a previous window/timestep are deleted.
\end{itemize}
It also writes \code{J\_value.txt} and \code{N\_obs.txt} (the matched-observation
count, used for the chi-square diagnostic). An \code{--obs-error-inflation}
factor multiplies all observation uncertainties, to account for correlated
retrieval errors.

\subsection{Gradient and update: \code{lbfgsb\_step\_monthly.py}}

Reads the adjoint accumulator \code{SurfaceFluxAdj\_CO2} at the month boundaries,
forms the monthly gradient by differencing (Section~\ref{sec:formats}), adds the
background gradient, performs one L-BFGS-B two-loop update with a unit step, and
writes the next state. Guards enforce correctness at runtime:
\begin{itemize}[nosep]
  \item each boundary record must exist at its exact timestamp (no nearest-record
        fallback, which could bias every monthly gradient by half a timestep);
  \item the accumulator must be $\approx 0$ at the window end (it is $0$ by
        construction);
  \item the monthly gradients must telescope to the whole-window gradient.
\end{itemize}
It also prints chi-square ($2J_{\mathrm{obs}}/N_{\mathrm{obs}}$) and
observation-vs-background gradient-balance diagnostics each iteration. The
\code{--n-months} option matches the control-vector length to the window.

\subsection{Diagnostics: \code{plot\_4dvar.py}}

Renders the state and convergence diagnostics. It infers the month count from the
state-vector size, so a monthly state is drawn as multi-panel (one map per month)
figures for the scaling field $\sigma$, the departure $\sigma-1$, and the
gradient, with a shared colour scale; panels are labelled by calendar month from
the window start. It also produces convergence histories, observation-minus-model
innovation maps and histograms, and animations across iterations. A plotting
failure is caught and never aborts the optimisation.

\subsection{Validation tooling}

\begin{itemize}[nosep]
  \item \code{check\_monthly\_sigma.py}: from a forward test run, verifies that
        the correct monthly field is applied to TER in the correct month, that no
        other flux component is affected, and that no wrong-month field leaks in.
  \item \code{fd\_validate\_gradient.run} + \code{fd\_gradient\_tools.py}: a
        finite-difference check that the differenced gradient (not the raw
        accumulator) matches a brute-force perturbation of $J$.
  \item \code{test\_monthly\_sigma.run}: the forward-only plumbing test that
        \code{check\_monthly\_sigma.py} evaluates.
  \item \code{adjoint\_determinism\_test.run}: reruns one adjoint twice from
        identical inputs to distinguish chaotic gradient noise from a
        nondeterminism bug.
\end{itemize}

% =============================================================================
\section{Per-iteration workflow}
\label{sec:workflow}
% =============================================================================

Each iteration runs the four phases in order. The last completed phase is written
to a progress file so that a killed or resubmitted job resumes at the right point.

\paragraph{Phase 1 --- sigma + forward.}
\code{write\_sigma\_monthly.py} writes the monthly fields from the current state.
The forward run directory is cleaned of stale checkpoints and outputs, the forward
GCHP integration runs over the window, and the emission-time adjoint checkpoints
it produces are moved (renamed, not copied) into the adjoint directory. The final
model state is moved to the adjoint restart.

\paragraph{Phase 2 --- forcing.}
Stale forcing files are deleted; all satellite-track files from \code{OutputDir}
(which was wiped before the forward run, so every track belongs to this iteration)
are passed to \code{co2\_adjoint\_forcing\_osse.py}, which writes $J_{\mathrm{obs}}$
and the per-timestep forcing. The observation-diagnostic archive is copied into the
per-iteration output directory.

\paragraph{Phase 3 --- adjoint.}
The adjoint run directory is cleaned, the backward GCHP integration runs (reading
the forcing and the moved checkpoints), and on success the consumed checkpoints are
deleted to bound disk usage. Deletion happens only \emph{after} the adjoint
succeeds, so a redo of this phase still has its inputs.

\paragraph{Phase 4 --- update.}
\code{lbfgsb\_step\_monthly.py} forms the gradient and the L-BFGS-B step and writes
the next state, which is snapshotted; plots are generated. Convergence
(gradient-norm or relative-cost tolerance) ends the run; otherwise the iteration
counter advances and the loop returns to Phase 1.

% =============================================================================
\section{On-disk data formats}
\label{sec:formats}
% =============================================================================

\paragraph{Monthly scaling fields.} \code{sigma\_monthly/YYYY/MM.nc4}, one
single-record file per month, variable \code{Sigma\_CO2}$(t,\mathrm{lat},
\mathrm{lon})$, \code{float32}, stamped at the month start.

\paragraph{State file.} A NumPy \code{.npz} with the L-BFGS-B state:
\code{x\_next} (next control vector), \code{x\_prev}, \code{g\_prev},
\code{J\_prev}, the limited-memory history \code{s\_hist}/\code{y\_hist},
\code{m\_used}, \code{iteration}, and the grid metadata
\code{nlat}/\code{nlon}/\code{nmon}/\code{start\_month}. The control vector is
$\code{nmon}\times\code{nlat}\times\code{nlon}$ in C order.

\paragraph{Adjoint output.} \code{GEOSChem.Adjoint.*.nc4} in the adjoint
\code{OutputDir}, containing \code{SurfaceFluxAdj\_CO2} on a $72\times46$ lat--lon
grid at the model timestep. The monthly gradient reads this at the thirteen month
boundaries $t_1,\dots,t_{12},t_{\mathrm{end}}$ and differences consecutive
boundaries:
%
\begin{equation}
  g_m(x) = A(x,t_m) - A(x,t_{m+1}),\qquad A(x,t_{\mathrm{end}})=0.
  \label{eq:gm}
\end{equation}

\paragraph{Forcing.} \code{CO2\_adjoint\_forcing\_*.nc4} in the forcing directory,
one per model timestep, read by the adjoint run via its input configuration; plus
\code{J\_value.txt} and \code{N\_obs.txt}.

% =============================================================================
\section{GCHP configuration and the forward/adjoint asymmetry}
% =============================================================================

The control variable is HEMCO scale factor \code{750} (\code{Sigma\_CO2}) applied
to the \code{TER\_CLASS\_CTEM} base entry. Of the active flux components, only TER
carries it. The forward and adjoint configurations differ in a way that is
essential to the gradient's correctness:

\begin{itemize}[nosep]
  \item \textbf{Forward} \code{HEMCO\_Config.rc}: the TER entry carries
        \code{750/751}, i.e.\ the monthly scaling (\code{750}) times the unit
        conversion (\code{751}). The forward \code{ExtData.rc} reads
        \code{Sigma\_CO2} from \code{sigma\_monthly/\%y4/\%m2.nc4} under a monthly
        refresh template, so the field changes at each month boundary.
  \item \textbf{Adjoint} \code{HEMCO\_Config.rc}: the TER entry carries \code{751}
        \emph{only}. The adjoint accumulator is weighted by
        \code{EmisCO2\_NetTerrExch} evaluated as the \emph{unscaled prior} TER flux
        at every timestep and every iteration. This is why the differenced gradient
        \eqref{eq:gm} needs no $1/\sigma_m$ factor, and why the adjoint never reads
        the monthly scaling.
\end{itemize}

\noindent\textbf{Corollary.} Only the forward control files require the monthly
plumbing; the adjoint configuration is unchanged from the single-window case
except that its output collection writes \code{SurfaceFluxAdj\_CO2} at the model
timestep (so a record exists at every month boundary) on the lat--lon grid.

% =============================================================================
\section{Monthly-specific plumbing and timing conventions}
% =============================================================================

\paragraph{Monthly input refresh.} The forward \code{ExtData.rc} entry for
\code{Sigma\_CO2} uses a monthly refresh template, reading a new file at each month
start. There is a one-timestep lag: the input component applies a new monthly file
one timestep after its nominal refresh stamp. On an instantaneous emissions
diagnostic this means the snapshot stamped exactly at a month boundary still
carries the previous month's scaling; the following record carries the new one.
This lag is benign (it is the same lag seen for other monthly inputs) but must be
accounted for when verifying the plumbing---the plumbing checker classifies the
boundary snapshot as belonging to the previous month for exactly this reason.

\paragraph{Month boundaries and the gradient.} With the emissions/output timesteps
at $30$ minutes, every month start is a multiple of the timestep, so an adjoint
record exists at each boundary and the exact-timestamp read in
\code{lbfgsb\_step\_monthly.py} always succeeds. The window ends on a month
boundary; because no integration step \emph{starts} at the window end, the input
component does not read a scaling file for the month after the window.

\paragraph{Multi-file forcing.} The forcing program consumes the monthly
satellite-track files separately and sums $J$; its month-coverage pre-checks are
the tripwire against stale or missing monthly tracks.

% =============================================================================
\section{Variable window parameterisation}
% =============================================================================

The month count is derived from the window rather than hard-coded. The driver
computes
\[
  \code{NMON} = (\text{end year}\times 12 + \text{end month})
              - (\text{start year}\times 12 + \text{start month}),
\]
and the start month from the window start, and passes \code{--n-months} (and
\code{--start-month}) to \code{write\_sigma\_monthly.py} and \code{--n-months} to
\code{lbfgsb\_step\_monthly.py}. The Python programs default to twelve months (the
full-year run) so existing callers are unaffected; the module-level constant is
retained for the finite-difference tooling that imports it. This is what allows a
short integration test (e.g.\ a three-month window) to exercise the entire
iteration at a fraction of the full-year cost. The window itself
(\code{T\_START}/\code{T\_END}) is overridable at submission and is propagated
across chained jobs.

% =============================================================================
\section{Running the system}
\label{sec:running}
% =============================================================================

The driver is submitted from \code{gchp\_4Dvar/}. Common invocations:

\begin{center}
\small
\begin{longtable}{@{}p{0.5\textwidth}p{0.44\textwidth}@{}}
\toprule
Command (\code{qsub \dots 4dvar\_optimizer.osse.monthly.run}) & Effect \\
\midrule
\endhead
(no options) & fresh full-year run, self-chaining \\
\code{-v RESTART=true} & resume at the recorded phase after a kill \\
\code{-v CHAIN=false} & do not self-resubmit when walltime is spent \\
\code{-v T\_END=2016-04-01,MAX\_ITER=1,CHAIN=false} & three-month, single-iteration integration test \\
\code{-v SIGMA\_B=0.2,MAX\_ITER=20} & set background error and iteration cap \\
\bottomrule
\end{longtable}
\end{center}

\noindent
Selected environment variables: \code{T\_START}/\code{T\_END} (window; must be
month boundaries), \code{NLAT}/\code{NLON} (control grid), \code{SIGMA\_B}
(background error), \code{MAX\_ITER} (global iteration cap across chained jobs),
\code{GTOL}/\code{FTOL} (convergence tolerances), \code{CHAIN} and \code{MAX\_CHAIN}
(self-chaining), \code{EST\_FORWARD\_S}/\code{EST\_ADJOINT\_S} (per-phase walltime
budget estimates).

% =============================================================================
\section{Operational notes}
% =============================================================================

\paragraph{Job chaining and walltime budgets.} A full-year iteration does not fit
in a single job's walltime, so the driver checkpoints phases and self-resubmits
with \code{RESTART=true} when the remaining walltime cannot accommodate the next
GCHP phase (gauged by \code{EST\_FORWARD\_S}/\code{EST\_ADJOINT\_S}). A killed job
redoes only the interrupted phase. \code{MAX\_CHAIN} caps runaway chains. When
running with \code{CHAIN=false}, the phase-budget estimates must be set below the
walltime, or the adjoint phase can be skipped: with chaining off, a failed budget
check exits cleanly rather than resubmitting, so a too-large estimate would end the
job after the forward without running the adjoint.

\paragraph{Checkpoint disk usage.} The forward run writes a large volume of
emission-time checkpoints for the adjoint (order a terabyte for a full-year run at
the fine timestep). The monthly driver deletes them after each successful adjoint,
so steady-state usage is bounded by one iteration in flight. The single-window
driver now does the same. Checkpoints are read by exact clock datestamp, so a run
at one timestep never reads checkpoints written at another---stale checkpoints from
a different timestep are inert but waste space and should be removed.

\paragraph{Isolation of variants.} The monthly and single-window optimisers use
separate run trees, control files, state files, forcing directories, logs, and
output directories. The finite-difference and plumbing-test jobs share the monthly
run tree with the optimiser and must not run concurrently with it; sequence them
with a job dependency.

\paragraph{Determinism.} The adjoint is bitwise-deterministic given identical
inputs; observed gradient noise is a genuine property of the convective adjoint,
not nondeterminism. This was established by rerunning an adjoint twice from
identical inputs.

% =============================================================================
\section{Cross-references}
% =============================================================================

\begin{itemize}[nosep]
  \item Scientific formulation, results, and outlook:
        \code{monthly\_4dvar\_science.tex}.
  \item Derivation and finite-difference verification of the monthly gradient
        \eqref{eq:gm}: \code{monthly\_sigma\_gradient.tex}.
\end{itemize}

\end{document}
